\clearpage
\scanpage{89}
\mspage{80}
% archive sheet number: 89

C'est le premier axiome de nature délicate, sur lequel nous \unsure{tombions|tombons}. Il est automatiquement satisfait pour les magmiens fidèles (car $\unsure{\operatorname{Omb}|\operatorname{Oub}}$ y est une fonction support ext. pleinement fidèle), plus \unsure{gén\textsuperscript{t}|génl\textsuperscript{t}} chaque fois qu'on a une fonction support externe pleinement fidèle. Le premier cas intéressant où il n'est pas sûr qu'on ait une telle fonction support, est celui du magmien associé à un ens ordonné, qu'il serait temps d'introduire ici en forme et d'examiner sous l'optique des supports.

Je voudrais donc à présent
\begin{enumerate}[label=\arabic*\textsuperscript{o}),leftmargin=*]
\item regarder le cas du magmien \deleted{\unsure{associé}} défini par un ens ordonné
\item Examiner le cas du magmien défini par $\leq$, $|\circ|$ seulement ($\overset{\circ}{\ll}$ étant défini via $\unsure{\leq|\ll}$, $|\circ|$, \deleted{\unsure{impliquant|indiquant} \unsure{une} \illegible{1}} \inserted{\unsure{entre des cinq.\ relations|entre diverses relations}} permettant de donner un sens à la relation $Y^{\circ}\subset X^{\circ}$ (en définissant $X^{\circ}=\unsure{\operatorname{supp}^{\circ}|\operatorname{Supp}^{\circ}}(X)\ldots$), donc de définir en conséquence $\overset{\circ}{\ll}$ \ldots
\item Voir si dans ces magmiens ``disjonctifs'', on n'aurait pas des exemples évidents où l'axiome des supports serait en défaut\ldots$^*$
\end{enumerate}
% slanted
\margin{$^*$ \uline{Oui} !}
